% ======================================================================
% Capítulo: Pruebas y validación
% Fases 7 y 8 del anteproyecto. Extensión orientativa: 4--6 páginas.
% ======================================================================
\chapter{Pruebas y validación}
\label{cap:pruebas}

Este capítulo describe cómo se comprueba que el sistema funciona y que cumple sus requisitos no funcionales: las pruebas automáticas de cada parte, las que protegen el aislamiento entre ayuntamientos y la privacidad, y las mediciones de rendimiento. La calidad de las respuestas se evalúa aparte, con RAGAS, en el capítulo~\ref{cap:estudio-chunking}.

\section{Estrategia de pruebas}
\label{sec:estrategia-pruebas}

Cada lenguaje del proyecto tiene sus pruebas. La integración continua ejecuta en cada cambio las de .NET, Go y Angular (sección~\ref{sec:calidad}), y un cambio no se despliega si falla cualquiera de ellas. Las del programa de evaluación en Python se ejecutan en local con \texttt{uv}, porque solo afectan a la evaluación automática del capítulo~\ref{cap:estudio-chunking}. La tabla~\ref{tab:pruebas} resume el conjunto.

\begin{table}[htbp]
\centering\small
\caption{Pruebas automáticas por proyecto}
\label{tab:pruebas}
\begin{tabular}{@{}l l r@{}}
\toprule
\textbf{Proyecto} & \textbf{Tipo} & \textbf{Pruebas} \\
\midrule
\texttt{Application.Tests} & Unitarias de los casos de uso & 304 \\
\texttt{ApiService.Tests} & Integración de la API, \emph{hubs} y seguridad & 150 \\
\texttt{Infrastructure.Tests} & Integración con PostgreSQL, Qdrant y RabbitMQ & 88 \\
\texttt{Domain.Tests} & Unitarias del dominio & 28 \\
\texttt{Worker.Tests} & Consumidores de ingesta & 27 \\
\texttt{ArchitectureTests} & Reglas de arquitectura & 16 \\
\texttt{Evaluation.Tests} & Programa de evaluación & 12 \\
\texttt{Gateway.Tests} & Pasarela (\emph{gateway}) & 5 \\
\midrule
\emph{Scraper} en Go & Unitarias, con detector de carreras & 47 funciones \\
\emph{Frontend} Angular & Componentes y servicios & 131 \\
Evaluación en Python & Validación del \emph{dataset} y puntuación & 16 \\
\bottomrule
\end{tabular}
\end{table}

\paragraph{Pruebas unitarias.} La mayor parte de las pruebas de .NET son unitarias, con xUnit, y se centran en la capa de aplicación: la orquestación del chat, las \emph{pipelines} RAG, la Arena, la ingesta y las estrategias de fragmentación. Las dependencias externas (modelos de lenguaje, almacén vectorial, mensajería) se sustituyen por dobles de prueba a través de las interfaces de la capa de aplicación, de modo que cada prueba comprueba la lógica sin red y en milisegundos.

\paragraph{Pruebas de integración.} Las pruebas que dependen del comportamiento real de la infraestructura se ejecutan contra contenedores efímeros creados con Testcontainers: la idempotencia de la ingesta en PostgreSQL y Qdrant, la concurrencia del estado de cada documento y de los votos de la Arena, el consumidor de RabbitMQ con sus reintentos y su cola de mensajes fallidos (DLQ), la migración del estudio de fragmentación y el informe de latencia operativa. Necesitan Docker; en un equipo sin él se omiten y en la integración continua se ejecutan siempre. Las pruebas de la API levantan el servicio completo en memoria con \texttt{WebApplicationFactory}, con autenticación simulada, y lo recorren por HTTP y por SignalR igual que un cliente real.

\paragraph{Pruebas de arquitectura.} Las pruebas de \texttt{ArchitectureTests} comprueban que el código respeta las reglas de la arquitectura limpia (sección~\ref{sec:clean}): que el dominio no depende de la persistencia, que la capa de aplicación no usa los SDK de la infraestructura, que cada proyecto solo referencia los paquetes permitidos, que la API no accede directamente a la base de datos ni fija \emph{prompts} de sistema y que el programa de evaluación no abre puertos. Estas pruebas también limitan el tamaño de los servicios del chat y su número de dependencias, para que la división hecha en la auditoría de código no se deshaga con el tiempo.

\paragraph{Go y Angular.} El \emph{scraper} se prueba con \texttt{go test -race}, que detecta accesos concurrentes sin sincronizar en el procesamiento paralelo de los documentos de un día, y con \texttt{go vet}. La integración continua comprueba además que el código generado a partir de los contratos gRPC está al día. El \emph{frontend} tiene pruebas de sus componentes y servicios, que se ejecutan sin navegador visible en cada cambio, y se compila en modo producción para detectar errores de tipos.

\section{Pruebas de seguridad y aislamiento}
\label{sec:pruebas-seguridad}

El requisito más delicado del sistema es que un ayuntamiento no pueda ver los datos de otro (sección~\ref{sec:seguridad}). Varias pruebas lo comprueban desde fuera, con peticiones reales a la API:

\begin{itemize}
  \item \textbf{Aislamiento por ayuntamiento.} Un usuario solo ve el historial de llamadas al modelo y la configuración de IA de su organización. Una sesión de chat no se puede abrir, leer, continuar ni borrar desde otro usuario, ni desde el mismo usuario en otra organización. Un voto sobre una batalla de la Arena de otra persona se responde como si no existiera y no la modifica.
  \item \textbf{Autenticación y autorización.} Las conexiones anónimas al chat se rechazan y las que no traen organización se prohíben fuera del entorno de desarrollo. Cada ruta de administración está clasificada como de plataforma o de organización; una prueba recorre todas y comprueba que un usuario anónimo recibe un 401, un usuario sin rol un 403 y que cada rol de administrador solo accede a sus rutas. Si alguien añade una ruta de administración sin clasificarla, la prueba falla.
  \item \textbf{Servicio interno.} La llamada gRPC del \emph{scraper} a la API para leer los filtros solo se acepta con la clave interna compartida.
  \item \textbf{Errores sin detalles.} Un identificador de sesión mal formado devuelve un error tipado, sin trazas ni detalles internos.
\end{itemize}

Dos pruebas de arquitectura protegen el aislamiento y la privacidad desde el código. La primera revisa que \texttt{IgnoreQueryFilters}, que salta el filtro por ayuntamiento, solo aparece en una lista de ficheros aprobados, cada uno con un filtro explícito o de uso exclusivo de los administradores de la plataforma. La segunda comprueba que la telemetría de los modelos de lenguaje no registra nunca el contenido de las preguntas y respuestas, solo los tokens y las duraciones: falla si alguna configuración, fichero de \emph{compose} o línea de código activa los diagnósticos sensibles. La pasarela tiene también sus pruebas: reenvía los errores del navegador al colector de telemetría sin las credenciales del usuario, limita las peticiones por cliente y rechaza los envíos demasiado grandes.

\section{Validación con usuarios reales}
\label{sec:validacion-usuarios}

El sistema está preparado para validarse con personal municipal: cada ayuntamiento es una organización con sus usuarios y sus datos separados, y la Arena permite comparar a ciegas las respuestas de dos \emph{pipelines} y votar la mejor (sección~\ref{sec:impl-arena}). Además, el \emph{shadow testing} genera comparaciones aunque los usuarios no voten, y la página de métricas muestra las victorias de cada \emph{pipeline} por periodo.

Esa validación no se ha llevado a cabo dentro de este trabajo. La primera versión se abrió durante dos días a algunas personas, de manera informal (sección~\ref{sec:scrum}), pero sin un protocolo ni datos que permitan sacar conclusiones. Una prueba con usuarios reales que diera resultados fiables exigiría muchas consultas y muchos votos durante semanas, y cada consulta en la Arena ejecuta dos \emph{pipelines}, una de ellas con varias llamadas al modelo. Con los modelos de un proveedor comercial, ese volumen supone un coste de API que no cabe en el presupuesto de un TFG: solo la ingesta inicial de los boletines con las \emph{pipelines} básica y jerárquica costó unos 30~USD entre OpenRouter y Google AI, sin contar las otras dos variantes ni las consultas (sección~\ref{sec:recursos}). Se canceló, por ese motivo, la apertura de la plataforma a los usuarios. Con modelos locales, la validación exigiría un servidor con GPU dedicado al servicio durante todo el periodo. Una validación pequeña, con modelos locales en el ordenador del capítulo~\ref{cap:estudio-chunking}, queda como la primera ampliación (sección~\ref{sec:futuras}).

Si se hiciera, la Arena debería ajustarse antes, porque hoy un voto no compara solo dos \emph{pipelines} (sección~\ref{sec:impl-arena}): habría que fijar el mismo modelo y el mismo \emph{prompt} en los dos lados, controlar la longitud y el formato de las respuestas, y agregar los votos con un modelo de comparaciones por pares, como el de Bradley y Terry \cite{bradley1952}, o con la tasa de victorias y su intervalo de confianza, exigiendo un número mínimo de votos antes de cambiar la \emph{pipeline} del chat. Mientras tanto, la comparación de las variantes se ha hecho con la evaluación automática del capítulo~\ref{cap:estudio-chunking}, cuyas métricas de recuperación no dependen de ningún juez.

\section{Rendimiento}
\label{sec:rendimiento}

Las cifras de rendimiento proceden de las ejecuciones del estudio de fragmentación, en el ordenador del capítulo~\ref{cap:estudio-chunking} con modelos locales, y se detallan allí (tablas~\ref{tab:resultados-latencia} y~\ref{tab:resultados-indexacion}). No son las de producción, que no tiene GPU y usa proveedores en la nube, ni incluyen la red ni SignalR, pero sí dan el coste relativo de cada \emph{pipeline} en condiciones iguales:
\begin{itemize}
  \item \textbf{Chat.} Con 5.994 documentos, las planas dan la respuesta completa en torno a un segundo en mediana y la jerárquica en unos 2,5\,s, porque reformula la consulta y responde más largo. Las planas cumplen con holgura el requisito de tiempo de respuesta (RNF12); la jerárquica cumple el de la respuesta completa, pero su primer texto llega hacia 1,03\,s, así que no cumple RNF12 tal como está redactado.
  \item \textbf{Ingesta.} Indexar un documento cuesta en mediana de 4\,s (tamaño fijo) a 19\,s (jerárquica), en segundo plano y sin afectar al chat. La colección completa se ingirió en las cuatro variantes sin perder mensajes.
  \item \textbf{Errores.} Ninguna de las 448 respuestas de evaluación falló. Para el uso real, el programa de evaluación genera un informe de latencia operativa a partir de los registros de la plataforma, con el que se podrán recalcular estas cifras cuando haya usuarios.
\end{itemize}
